\section{Related Work}

\subsection{Onboard science autonomy and rover target selection}
Onboard science autonomy systems such as AEGIS-style target selection aim to reduce reliance
on delayed ground loops by identifying and prioritizing scientifically interesting
targets~\citep{estlin2012,francis2017aegis}.
Operational ChemCam targeting on Curiosity demonstrated that onboard selection can be adopted
as a routine science tool~\citep{francis2017aegis}, and analogue-mission studies have examined
the utility of rover science-autonomy capabilities under reduced ground interaction~\citep{francis2019utility}.
Human-in-the-loop operations remain essential for mission-level decisions, but bandwidth and
latency constraints motivate onboard screening that is explicit about uncertainty and
false-positive control.
ARTPS follows this operational framing: anomaly cues feed prioritization rather than replacing
mission planning.

\subsection{Visual anomaly detection}
Feature-distribution and memory-bank methods such as PaDiM and PatchCore are widely used
industrial anomaly baselines~\citep{defard2020padim,roth2022patchcore}.
Reconstruction- and embedding-based detectors similarly score unusual appearance, but a high
anomaly score is not identical to an operationally actionable science target under rover
constraints.
ARTPS therefore treats anomaly responses as cues inside a broader prioritization stack that
includes fusion, localization, suppression, diversity, and buffering.

\subsection{Monocular relative depth}
Dense monocular depth models provide image-relative structure useful for near/far
organization~\citep{ranftl2021dpt}.
They do not, by themselves, provide metric distance, physical object size, or calibrated 3D
reconstruction under planetary scale ambiguity.
ARTPS uses relative depth as an ordering and policy cue; no metric-distance claim is made.

\subsection{Edge onboard AI and safety-aware autonomy}
Onboard compute, thermal, and power limits constrain which models can run continuously.
Self-reliant rover concepts further emphasize maintaining productivity with reduced
ground-in-the-loop interaction~\citep{gaines2020srr}.
Uncertainty, fail-safe ``no selection'' behaviour, and runtime monitoring matter as much as
ranking discrimination for operational use.
ARTPS therefore states only bounded assurance within a defined operational envelope,
not universal or flight-certified guarantees.
